%% APPENDIX D -- the false-alarm expectation and its resampling, block
%% quality, registration, the pairs and the drift grid's edges.  Carries
%% app:falsealarm (04_method.tex, 05_results.tex, 08_backmatter.tex,
%% app_C, app_I), eq:expect, eq:pnull,
%% sec:blockquality (04_method.tex, 05_results.tex),
%% sec:regbudget (05_results.tex), sec:pairacct (cited by v4.30 Sec. 4; kept
%% so the label resolves either way) and sec:edge (04_method.tex, 05_results.tex,
%% pairs_v415.py).
%% r30 (write-apx): D.1 about half its v4.30 length with the estimator and
%% the resampling as two displayed equations; D.2 and D.3 short; the pair
%% accounting kept short here because Sec. 4 cites it; Table D1 (tab:edge) is
%% one sentence of D.5 (r34: D.5 folded into D.4; sec:pairacct kept as a phantom label).  r34 (W3): D.1 split one idea per paragraph; the
%% rule that chose the retained control positions and the stratified share
%% (R2-5); new D.2 sec:choicegrid + tab:choicegrid (R2-2, 8-cell grid).
%% r35 (W3, A1): the per-window attributed share, drawn towards the pooled
%% share (beta-binomial, strength by maximum likelihood), is the ADOPTED
%% estimator (\CeExp, \CeProbNull); the pooled share (\CeExpPool,
%% \CeProbNullPool) and the median-split strata are the two checks, in D.1
%% paragraph 3; \CeMaskPct is the share of the expectation removed with the
%% per-window shares (no longer 89/286 = \CeAwPool).  \CgAwShr* are retired.
%% The control assumption is stated once, in
%% Appendix B.  The band-edge sentence is REQUIRED on the page (v409 V7c).
\section{The false-alarm expectation and the control ensemble}
\label{app:falsealarm}

\subsection{The expectation and its resampling}

The comparison with chance asks whether the stars cross the trigger more
often than empty sky in their own fields. Each window supplies
its own answer to how often empty sky does so, from its control positions,
and one attribution set serves both sides, the mask of
Appendix~\ref{app:linemask} being applied to a control maximum exactly as to
a crossing. Over the \CeNWin{} Class~A windows $w$ of the primary census the
expected number of unattributed crossings is
\begin{equation*}
\bar N=\tau\sum_{w}r_w\,(1-a_w),\qquad r_w=\frac{n_w}{N_{\rm ctrl}},
\end{equation*}
where $n_w$ of the window's $N_{\rm ctrl}=\NCtrl$ control positions reach the
trigger, $a_w$ is the window's share of control exceedances that land
inside the mask, and $\tau$ is the share the atmospheric screen leaves. The rates alone
give \CeExpTrig{} expected crossings.

The attributed share has to be measured where control maxima actually fall
in frequency, and that needs their spectra. The extractor kept a dynamic
spectrum for only \RcNRetCtrl{} of the control positions in each window, to
save storage, and it chose them when the products were written, as entries
spaced evenly through its list of control positions. That list is in the
order of the seeded random draw that placed the positions, so the choice
depends on nothing in the data. Of the \CeNExc{} exceedances at these
positions, \CeNExcAttr{} lie inside the mask.

That share is not the same in every window. Windows differ in how much of
their band the mask covers, and a window crowded with listed lines loses
more of its control maxima to the mask than one that is not. Each window's
$a_w$ is therefore measured from its own exceedances. Most windows have only
a few, so each window's share is drawn towards the pooled share, the more
strongly the fewer exceedances it has, and the strength of that pull is the
one that best fits the scatter between windows (a beta-binomial model).
Taken window by window in this way, attribution removes \CeMaskPct{} per
cent of the expectation, and this is the expectation used throughout.

Control spectra were kept in \CeNExcWin{} of the \CeNWin{} windows. The
rest carry the share of their own channels that the mask covers, whichever
reading is used.

Two simpler readings are kept as checks. One pooled share, \CeAwPool, for
every window whose control spectra were kept gives \CeExpPool{} expected,
with $P=\CeProbNullPool$. Shares measured
separately for the windows above and below the median masked fraction,
\CgAwHi{} and \CgAwLo{} from \CgAwNExcHi{} and \CgAwNExcLo{} exceedances,
give \CgAwStratExp, with $P=\CgAwStratProb$, \CgAwStratClause. Both checks
leave the count consistent with chance.

The screen (Appendix~\ref{app:rfi}) then removes \CeTlCtrlCut{} of the
\CeTlCtrlSurv{} exceedances that remain at these positions, leaving $\bar N=\CeExp$ against
\CeObs{} observed.

Over all windows together, the exceedance rate near a band edge is
$\times\VnEdgeRatio$ the interior rate, but one window, toward HD~48370, holds \VnEdgeDomPct{} per
cent of the near-edge exceedances. Without it the ratio is
$\times\VnEdgeNoDomRatio$ ($\times\VnEdgeHalfCleanRatio$ on a held-out half,
$\pv{\VnEdgeHalfCleanP}$), so no edge correction is applied.

Some data enter neither side of the comparison. The \CeNWinTail{} Class~A
windows of the archival tail are left out, and so is the crossing toward
\CtUncatStar. That crossing is \MbUncatAttrWord{} and is one of the \CsNCross{}
census crossings, but its position has no catalogue row. It therefore has no
control positions of its own on the expected side, and so it is not counted
on the observed side either. The window it lies in is still summed once,
through the control positions about \MbUncatHost. Which statistic every
crossing enters is listed in \texttt{statistics/\allowbreak
event\_membership.csv} of the data archive.

Windows share execution blocks, so the count is compared through block
resampling. In each of $B=\CeNBoot$ resamplings the blocks are drawn with
replacement, each window's rate is redrawn binomially from its own probes,
and the attributed shares and the screen are redrawn from their own counts,
giving $\bar N_b$; then
\begin{equation*}
P=\frac{1}{B}\sum_{b=1}^{B}\Pr\!\left(N\ge N_{\rm obs}\mid\bar N_b\right),
\end{equation*}
with $N$ a Poisson count and $N_{\rm obs}=\CeObs$, gives $P=\CeProbNull$,
conditional on the control-position model (Appendix~\ref{app:radius}). As
a check, a pseudo-star read at one control position in every window of
a block is no more dispersed than one read window by window (variance ratio
\BgDisp) and reaches \CeObs{} with probability \BgPseudoP. The widened recurrence window
of \S\ref{sec:etacrv} is read in the same way over its \WuNWin{} distinct
windows.

The bound quoted in \S\ref{sec:trigunif} reads the lower tail of the same
average over resamplings. With $k$ added to the mean $\bar N_b$ of every resampling,
\begin{equation*}
P_{\le}(k)=\frac{1}{B}\sum_{b=1}^{B}\Pr\!\left(N\le N_{\rm obs}\mid\bar N_b+k\right)
\end{equation*}
is the probability, on the same control positions, of a count as low as the
one observed, and it falls as $k$ grows while $P$ rises. It is
\CeProbAtLim{} at $k=\CeAddLim$ and drops to \CeProbAboveLim{} for one more.
So \CeAddLim{} is the largest additional mean number of primary crossings
compatible with the observed count at the one-sided \CeAlphaPct{} per cent
level, conditional on the mask and on the control-position model. It is not
a confidence bound on a realised number of crossings. Both tails are
tabulated against $k$ in
\texttt{statistics/\allowbreak chance\_allowance\_table.csv}.

The control positions of one window are not independent draws, since
neighbouring positions see much the same noise. Their layout can be rebuilt
exactly from the extractor's seeded random positions, and two positions
less than a tenth of a synthesised beam apart are then found to be
correlated at \CcRhoNear. The correlation falls to half that by
\CcHalfBeam\,$\theta_{\rm syn}$, and beyond half a beam no binned value
exceeds \CcRhoHalfBeamOut. The correlogram is given in
\texttt{statistics/\allowbreak control\_correlation.csv}.

To see what this correlation does to the comparison, each window's positions
were also grouped into blocks, either square cells from half a beam to a few
beams across or arcs of the annulus. The \CcNLead{} windows that carry
\CcLeadPct{} per cent of $\bar N$ then keep a median of \CcNeffLeadMed{}
effective positions out of \NCtrl, and never fewer than \CcNeffLeadMin.
Drawing each window's rate from its own blocks instead of binomially raises
the part of the scatter in $\bar N$ due to the window rates from
\CcSdBinom{} to \CcSdCorr. That is small beside the \CcSdTotal{} of the block
resampling as a whole. With the correlation between control positions kept,
$P$ is \CcProbRange{} and the bound \CcAddLimRange, close to the \CeProbNull{}
and \CeAddLim{} of the binomial redraw, which is therefore kept.
Each window's effective number of positions is given in
\texttt{statistics/\allowbreak control\_correlation\_windows.csv}.

\subsection{The choices made after the crossings were seen}
\label{sec:choicegrid}

Two choices behind the comparison were settled after the crossings had been
examined, and each could have been made the other way. The control field
can be the whole primary beam or only the positions nearest the star, and
the mask can include SO or leave it out. The comparison was repeated for
each of the \CgNTrialWord{} combinations of field and mask, with the same
resampling. Each combination was also read with the expected rate raised by
the star's excess just below the trigger (Appendix~\ref{app:radius}). That
excess is a measured factor and was not chosen, so it is a check on the
background. Table~\ref{tab:choicegrid}
gives all \CgNCellWord{} readings.

The smallest probability in the table is \CgMinProb, for \CgMinCell.
Multiplied by \CgNTrialWord, a Bonferroni factor that allows for its having
been picked out of the table, it is \CgBonfPhrase, \CgVerdictClause.
Only the field and the mask count towards that factor. The stellar rate
is one measured factor, \CeNear, the same in every combination, and raising
the expected count can only raise the probability. The smallest value therefore always lies among the
\CgNTrialWord{} combinations read at the control rate.

\begin{table}
\centering
\scriptsize
\setlength{\tabcolsep}{3pt}
\caption{Unattributed crossings observed and expected, with the chance
probability, for each control field and line mask.}
\label{tab:choicegrid}
% GENERATED by chancegrid_r34.py (round 3420) -- do not hand-edit.
\begin{tabular}{@{}llrrrrr@{}}
\hline
field & mask & $N_{\rm obs}$ & $\bar N$ & $P$ & $\bar N\times\CeNear$ & $P$ \\
\hline
whole beam & with SO (adopted) & \CeObs{} & \CeExp{} & \CeProbNull{} & \CeExpNear{} & \CeProbNear{} \\
whole beam & without SO & \MkNoSOObs{} & \MkNoSOExp{} & \MkNoSOProb{} & \CgExpWOSt{} & \CgProbWOSt{} \\
nearest field & with SO & \FrNarrowObs{} & \FrNarrowExp{} & \FrNarrowProb{} & \CgExpNSSt{} & \CgProbNSSt{} \\
nearest field & without SO & \CgObsNO{} & \CgExpNO{} & \CgMinProb{} & \CgExpNOSt{} & \CgProbNOSt{} \\
\hline
\end{tabular}

\end{table}

\subsection{Execution-block quality}
\label{sec:blockquality}

The comparison assumes that each block's control ensemble behaves like
noise, and in one block all \NCtrl{} probes exceed the trigger. Every block
is therefore judged on the median of its windows' control ensembles over its
class median. The same measure, applied window by window, also flags
\DqNWin{} of the \NWindows{} windows at \BqFactor{} the survey median. That is
a flag on single windows inside blocks that pass, and it is why the crossing
toward \BkFlagOtherStar{} is in quality-passing data and yet is withheld from
the tests that depend on amplitude (\S\ref{sec:technosearch}). The criterion was written after the condition
had been seen, but the choice is not a delicate one. The failing block sits
at \BqFailQ{} and the next at \BqNextQ, and its \BqFailNWin{} coarse-channel
windows hold no census crossing once ALMA's final flags are restored
(\S\ref{sec:method}).

\subsection{Registration}
\label{sec:regbudget}

Between epochs, the predicted cell is registered to \RgBudgetChMed{} of a
channel in every window. Widening each reading to the worst cell within
\RgTolKms\,km\,s$^{-1}$, re-anchoring it along the discovery track, or
raising the bar to \ZtTrialThr{} changes no disposition, although
registration to a single channel is tested only on the unresolved line of
\RgRegStar, which registers to \RgRegMaxCh{} of a channel.

\phantomsection\label{sec:pairacct}%
Where the epochs of a pair of crossings (\S\ref{sec:etacrv}) are close
enough for the predicted offset to be tighter than the searched window, all
\TjNPower{} such pairs miss a drift changing at a constant rate by at least
\TjSigMin{} times the resolution.

\subsection{The drift grid's edges}
\label{sec:edge}

The drift trials run from $-\dot\nu_{\max}$ in half-channel steps to the
first trial at or beyond $+\dot\nu_{\max}$, so the searched interval is
asymmetric by up to one step. Of the \EgNPop{} crossings of \EgDriftPop,
\EgNDrift{} have a fitted drift, and of these the \EgNEdge{} toward
\EgEdgeStars{} lie on an outermost trial of their own window's grid. Their
drift rates, in the observatory frame, and the accelerations they imply are
lower bounds, and \EgEdgeOutClause.
